\section{Pipeline Overview}

All trained architectures use the same split, preprocessing objects, optimizer family, trainer, checkpoint rule, and test evaluator. Figure~\ref{fig:pipeline-overview} shows the complete execution path. Training batches are augmented and normalized; validation and test batches are normalized without augmentation. At every epoch, the trainer records loss and accuracy, saves a checkpoint only when validation loss improves by at least $10^{-4}$, and stops after ten consecutive non-improving epochs. The best-validation-loss checkpoint is restored before the official test set is evaluated.

\begin{figure}[H]
    \centering
    \includegraphics[width=\textwidth,height=0.72\textheight,keepaspectratio]{graphics/pipeline-overview.pdf}
    \caption{End-to-end Fashion-MNIST experiment pipeline used by every completed architecture.}
    \label{fig:pipeline-overview}
\end{figure}

The implementation separates responsibilities across \texttt{dataloader.py} (splits and transforms), \texttt{models.py} (architectures), \texttt{trainer.py} (training, validation, early stopping, and checkpoints), \texttt{metrics.py} (predictive and resource metrics), and \texttt{plotter.py} (EDA and result figures). This separation allows the model family to change without changing the evaluation protocol.
